\documentclass[12pt]{article}
\usepackage{fontspec}
\usepackage{polyglossia}
\setmainlanguage{french}
\usepackage{csquotes}
\usepackage{lipsum}
\usepackage{array}% nécessaire pour les colonnes de largeur fixée avec alignement : >{...}p{...}

\begin{document}

\section{Insérer un tableau}

\lipsum[1]

% 1-2. Le code produit par l'assistant de TeXstudio :
%    - {|l|c|r|} : trois colonnes (gauche, centre, droite), séparées par des traits verticaux ;
%    - & : passage à la colonne suivante ; \\ : fin de ligne ; \hline : trait horizontal ;
%    - \multicolumn{2}{|c|}{texte} : fusionne deux colonnes, centrées, avec un trait de chaque côté.
%    Ajouter une ligne : recopier une ligne existante (même nombre de &), suivie de \hline.
% 4. Environnement table : le tableau devient un flottant, avec une légende (\caption).

\begin{table}[ht]
	\centering
	\begin{tabular}{|l|c|r|}
		\hline
		\multicolumn{2}{|c|}{Lieu et date} & Événement \\
		\hline
		Paris & 1789 & Prise de la Bastille \\
		\hline
		Varennes & 1791 & Arrestation du roi \\
		\hline
		Valmy & 1792 & Victoire française \\% ligne ajoutée à la main
		\hline
	\end{tabular}
	\caption{Quelques événements de la Révolution}
\end{table}

\lipsum[2]

% 3. Colonne de largeur fixée : p{4cm} (le texte revient à la ligne dans la colonne).
%    L'assistant ajoute >{\raggedright\arraybackslash} pour aligner ce texte à gauche (package array).
%    Pour modifier la largeur, il suffit de changer la valeur : p{4cm} -> p{6cm}.

\begin{table}[ht]
	\centering
	\begin{tabular}{|c|>{\raggedright\arraybackslash}p{6cm}|}
		\hline
		Date & Événement \\
		\hline
		1789 & Prise de la Bastille, le 14 juillet, après l'émeute du faubourg Saint-Antoine \\
		\hline
		1792 & Proclamation de la République, le 21 septembre \\
		\hline
	\end{tabular}
	\caption{Un tableau avec une colonne de largeur fixée}
\end{table}

% 5. Liste des tableaux (deux compilations sont nécessaires).
\listoftables

\end{document}
